\section{Repository Blueprint, Milestones, and Release Gates}
\label{app:repository}

This appendix is a build specification for an implementation agent. The tree is normative in responsibility, not necessarily in exact filenames.

\subsection{Repository tree}

The following decomposition assigns one enforceable responsibility to each package and keeps semantic IR, planning, execution, storage, backends, and operator plugins independently testable.

\begin{lstlisting}
ams/
  pyproject.toml
  CMakeLists.txt
  Cargo.toml                         # if Rust native core is selected
  LICENSE
  SECURITY.md
  CONTRIBUTING.md
  CODE_OF_CONDUCT.md
  README.md
  schemas/
    manifest.schema.json
    plan.schema.json
    config.schema.json
    trace.schema.json
  src/ams/
    __init__.py
    api.py
    config.py
    errors.py
    version.py
    ir/
      graph.py
      values.py
      effects.py
      guards.py
      serialization.py
      passes/
        canonicalize.py
        aliases.py
        shapes.py
        liveness.py
        fusion.py
    compiler/
      capture_pytorch.py
      import_checkpoint.py
      pattern_registry.py
      lower.py
      candidate.py
      planner.py
      cost_model.py
      plan_cache.py
    runtime/
      engine.py
      session.py
      executor.py
      scheduler.py
      memory.py
      leases.py
      events.py
      virtual_tensor.py
      transactions.py
      kv_manager.py
      cancellation.py
    storage/
      base.py
      manifest.py
      ams_container.py
      safetensors_store.py
      mmap_store.py
      file_store.py
      direct_store.py
      cache.py
      checksum.py
      journal.py
    backends/
      base.py
      cpu.py
      cuda.py
      rocm.py
      metal.py
      registry.py
    ops/
      base.py
      linear.py
      embedding.py
      normalization.py
      attention.py
      mlp.py
      moe.py
      convolution.py
      ssm.py
      reductions.py
      indexing.py
      sampling.py
    codecs/
      base.py
      raw.py
      fp_cast.py
      int8.py
      int4.py
      lossless.py
    training/
      autograd.py
      rematerialization.py
      gradients.py
      optimizers.py
      checkpoint.py
    integrations/
      pytorch.py
      transformers.py
      generation.py
    telemetry/
      metrics.py
      trace.py
      logging.py
      analysis.py
    cli/
      main.py
      convert.py
      inspect.py
      preflight.py
      run.py
      benchmark.py
      trace.py
      recover.py
      doctor.py
  native/
    include/ams/
      abi.h
      status.h
      backend.h
      storage.h
      scheduler.h
      memory.h
    src/
      core/
      storage/
      cpu/
      cuda/
      rocm/
      metal/
    tests/
  tests/
    unit/
    property/
    numerical/
    integration/
    models/
    memory/
    stress/
    fault/
    recovery/
    security/
    conformance/
  benchmarks/
    micro/
    operators/
    models/
    analysis/
    baselines/
  fuzz/
    manifest/
    plan/
    codecs/
    abi/
  docs/
    architecture/
    adr/
    api/
    operators/
    backends/
    formats/
    security/
    operations/
    performance/
    model_support/
  examples/
    cpu_out_of_core/
    gpu_subtensor/
    kv_storage_paging/
    custom_operator/
  tools/
    generate_schemas.py
    check_resource_formulas.py
    build_repro_bundle.py
  .github/workflows/
    lint.yml
    unit.yml
    native-sanitizers.yml
    cpu-integration.yml
    cuda-integration.yml
    rocm-integration.yml
    fuzz.yml
    docs.yml
    package.yml
    benchmark-regression.yml
    release.yml
\end{lstlisting}

\subsection{Dependency rules}

\begin{enumerate}
  \item \code{ir} has no dependency on framework, storage, or backend implementations.
  \item \code{compiler} depends on IR and abstract capability/storage metadata, not runtime instances.
  \item \code{runtime} depends on abstract backends/stores and serialized plans, not PyTorch graph objects.
  \item \code{ops} plugins depend on IR/operator contracts and backend kernel interfaces.
  \item framework integrations depend on public API/compiler import surfaces and may be optional packages.
  \item native core exposes no Python callback on device completion paths that can block while holding broker locks.
  \item tests may depend on internal modules; examples and third-party plugins use only public APIs.
\end{enumerate}

A static architecture test checks imports and native library dependencies.

\subsection{Implementation milestones}

\paragraph{Milestone 0: specification executable skeleton.}
Set up schemas, typed errors, virtual tensor descriptors, plan serialization, CPU-only CI, docs, and a synchronous reference executor. No accelerator claim. Acceptance: exact scalar linear, reduction, elementwise, and storage round-trip under a tiny bounded arena.

\paragraph{Milestone 1: CPU external-memory core.}
Implement range store, AMS container, mmap/file adapters, broker/leases, synchronous then asynchronous scheduler, CPU vector kernels, converter, and invariant tests where weights/input/output exceed the DRAM arena. Acceptance: a small transformer graph executes from storage with exact reference parity.

\paragraph{Milestone 2: two-dimensional linear accelerator path.}
Implement CUDA first or the selected primary accelerator: arena, events, pinned ring, async copies, device linear kernels, buffer pipeline, and memory accounting. Acceptance: a single weight row and full output can each exceed device arena; no whole-tensor allocation occurs; reference parity passes.

\paragraph{Milestone 3: complete dense transformer inference.}
Add embedding, norms, RoPE, gated MLP, exact attention, virtual KV, LM-head streaming, sampling, graph compiler, and Transformers adapter. Acceptance: supported decoder-only models execute with checkpoint larger than VRAM and, on a configured test, larger than DRAM cache.

\paragraph{Milestone 4: performance planner.}
Add calibration, candidate Pareto sets, graph fusion, prefill/decode plans, caching, continuous batching, direct IO optional path, trace analysis, and baselines. Acceptance: no unexplained pipeline serialization; plan estimates and measured stage times meet published tolerance.

\paragraph{Milestone 5: model-family breadth.}
Add MoE, convolution, SSM/recurrent, grouped/batched operators, dynamic-shape plan cache, second accelerator backend, and custom-op SDK. Acceptance: every advertised model has an operator coverage report and conformance artifact.

\paragraph{Milestone 6: training.}
Add joint forward/backward lowering, activation spill/rematerialization, gradient tiles, optimizer streaming, transactional checkpoints, and distributed composition. Acceptance: restart-exact small-model training and gradient/optimizer differential tests.

\paragraph{Milestone 7: enterprise hardening.}
Complete multi-tenant quotas, signatures, sandboxed converter, fuzzing, recovery tools, signed packages, SBOM, upgrade tests, long-duration stress, and incident runbooks.

Milestones are cumulative. A later feature cannot weaken the memory invariant of earlier guaranteed paths.

\subsection{Coding standards}

Python uses strict type checking, formatting/linting, immutable descriptors, explicit exceptions, and no implicit global singleton backend. Native code uses RAII/ownership types, checked arithmetic, sanitizers, warning-as-error, and zero undefined behavior in parser/scheduler paths. Kernels document shape/alignment/workspace contracts and have scalar references.

No TODO, placeholder pass, silent \code{except}, or ``temporary'' unbounded allocation is permitted in a release path. Experimental features are feature-gated and excluded from the support matrix.

\subsection{Test and CI matrix}

Every pull request runs:

\begin{itemize}
  \item formatting, lint, type, schema, documentation, license, and architecture checks;
  \item CPU unit, property, numerical, memory, cancellation, and parser tests;
  \item native ASan/UBSan/TSan where compatible, plus leak checks;
  \item deterministic reproducibility and plan serialization tests;
  \item a small end-to-end CPU out-of-core example.
\end{itemize}

Scheduled/self-hosted CI runs CUDA and ROCm conformance, direct IO, multi-device, long stress, fault/recovery, fuzz corpora, and benchmark regression. Release CI rebuilds from a signed tag, verifies reproducibility metadata, generates SBOM/provenance, signs artifacts, and executes smoke installs in clean environments.

\subsection{Regression policy}

Microbenchmarks have noise-qualified baselines by hardware pool. A regression gate uses robust statistics and requires investigation beyond a configured threshold. Baseline updates include a rationale, before/after traces, and confirmation that memory/correctness did not regress. Performance improvements that violate resource bounds are rejected.

\subsection{Definition of done by component}

A component is done only when it has:

\begin{enumerate}
  \item public or internal contract and invariants;
  \item implementation with bounded resources;
  \item unit, property, differential, cancellation, and fault tests as applicable;
  \item metrics and structured diagnostics;
  \item documentation and example;
  \item compatibility/versioning declaration;
  \item benchmark or cost-model calibration when on a critical path;
  \item security review for parsers, native code, storage, or plugins.
\end{enumerate}

\subsection{First production slice}

The smallest defensible production slice is not the original row-chunk module. It is:

\begin{itemize}
  \item CPU reference plus one accelerator backend;
  \item AMS/safetensors range storage and streaming converter;
  \item brokered arenas, events, scheduler, and traces;
  \item true 2D linear, embeddings, RMSNorm/LayerNorm, RoPE, gated MLP, exact attention, virtual KV, LM head, greedy/top-$k$ sampling;
  \item functional graph import for one well-defined decoder family;
  \item model/row/output/KV oversize conformance tests;
  \item preflight with exact bounds and CPU fallback.
\end{itemize}

Anything less can be a prototype, but should not claim the core invariant for arbitrary supported LLMs.
